Consistent with the preceding per-electrode findings, the region-collapsed results revealed a pronounced anterior-to-posterior decline in detection performance, as shown in Table~\ref{tab:region_performance}. For Internal, the mean $F_1$ decreased from 69.68\% across the six frontal electrodes to 14.80\% across the six central electrodes, 1.94\% across the four parietal electrodes, and 2.30\% across the eight occipital electrodes. Cao2018 exhibited a similar monotonic decline, with the mean $F_1$ decreasing from 68.91\% for the six frontal electrodes to 38.13\% for the six central electrodes, 18.70\% for the two parietal electrodes, and 8.24\% for the six occipital electrodes.

\begin{figure}[ht]
    \centering
    \includegraphics[width=0.9\linewidth]{fig_region_performance.pdf}
    \caption{Per-electrode detection $F_1$ under the full montage, shown as two stacked bar
        charts, Internal on top and Cao2018 on the bottom, with one bar per electrode of the
        32-channel array, coloured by scalp region and labelled with its $F_1$ value. Each
        bar is the macro $F_1$ for that electrode, scored inside the full-montage run and
        averaged across sessions.}
    \label{fig:region_performance}
\end{figure}
